\documentclass[12pt,a4paper]{article}
\usepackage[utf8]{inputenc}
\usepackage[T1]{fontenc}
\usepackage[brazil]{babel}
\usepackage{graphicx}
\usepackage{tabularx}
\usepackage{listings}
\usepackage{hyperref}
\usepackage{geometry}
\geometry{margin=2.5cm}

\title{Relatorio Tecnico \\ CPU RISC-V RV32I em Verilog}
\author{Projeto EEL580 -- Arquitetura de Computadores}
\date{}

\begin{document}
\maketitle
\tableofcontents
\newpage

\section{Introducao}

Este projeto descreve a implementacao de uma CPU RISC-V de 32 bits (RV32I) com
pipeline de 5 estagios, utilizando Verilog para integracao no simulador Digital,
atendendo aos requisitos do documento \texttt{Projeto\_CPU\_p1.pdf}.

\section{Premissas Adotadas}

\begin{itemize}
    \item O subconjunto de instrucoes RV32I exigido pelo PDF foi implementado:
          \texttt{add}, \texttt{addi}, \texttt{auipc}, \texttt{sub},
          \texttt{and}, \texttt{andi}, \texttt{or}, \texttt{ori},
          \texttt{xor}, \texttt{xori}, \texttt{sll}, \texttt{slli},
          \texttt{srl}, \texttt{srli}, \texttt{lw}, \texttt{lui},
          \texttt{sw}, \texttt{jal}, \texttt{jalr}, \texttt{beq}, \texttt{bne}.
    \item As memorias de instrucao e dados sao externas a CPU, expostas via
          interfaces separadas (\texttt{imem\_*} e \texttt{dmem\_*}),
          permitindo o uso das memorias ROM e RAM do Digital.
    \item O sinal \texttt{load\_enable\_i} congela toda a CPU durante a carga
          assincrona das memorias, preservando o estado interno.
    \item O reset e assincrono (\texttt{reset\_i}), zerando PC, pipeline e
          banco de registradores.
    \item Nao ha modo supervisor (S Mode); apenas o pipeline de inteiros.
    \item A ALU foi implementada em Verilog, sem uso de blocos combinacionais
          multi-bit do Digital, conforme recomendado.
\end{itemize}

\section{Arquitetura do Pipeline}

A CPU foi organizada em 5 estagios: IF, ID, EX, MEM e WB.

\subsection{IF (Instruction Fetch)}
O \texttt{program\_counter} mantem o PC atual e calcula PC+4. O endereco da
instrucao e apresentado em \texttt{imem\_addr\_o} e a instrucao e lida de
\texttt{imem\_data\_i}. Em caso de branch tomado ou jump, o PC e redirecionado
para \texttt{target\_addr\_i}.

\subsection{ID (Instruction Decode)}
O \texttt{instruction\_decoder} extrai opcode, rd, rs1, rs2, funct3, funct7 e
gera o imediato sign-extended para os formatos R, I, S, B, U e J. A
\texttt{control\_unit} decodifica o opcode e gera os sinais de controle. O
\texttt{register\_file} realiza leitura assincrona de rs1 e rs2, com x0 sempre
retornando zero. Um mux de internal forwarding WB$\rightarrow$ID seleciona o
dado do WB quando ha conflito write-before-read.

\subsection{EX (Execute)}
A \texttt{forwarding\_unit} seleciona a origem dos operandos A e B da ALU:
valor do banco (00), forward do MEM (01) ou forward do WB (10). Para
\texttt{auipc}, a entrada A recebe o PC; para \texttt{lui}, recebe zero. A
\texttt{alu} executa ADD, SUB, AND, OR, XOR, SLL, SRL e PASS\_B. O
\texttt{branch\_comparator} compara rs1 e rs2 para \texttt{beq} e
\texttt{bne}. O endereco de destino e calculado como PC+imm (branches/JAL)
ou rs1+imm (JALR, com bit 0 = 0).

\subsection{MEM (Memory Access)}
O endereco de memoria de dados e o resultado da ALU. Para \texttt{sw}, o dado
de rs2 e escrito. Para \texttt{lw}, o dado e lido. A escrita e suprimida quando
\texttt{load\_enable\_i} esta ativo.

\subsection{WB (Write Back)}
Para \texttt{jal}/\texttt{jalr}, o valor escrito e PC+4 (endereco de retorno).
Para \texttt{lw}, o dado lido da memoria. Para demais instrucoes, o resultado
da ALU.

\section{Tratamento de Hazards}

\subsection{Data Hazards (RAW)}
\textbf{Forwarding EX/MEM e EX/WB}: A \texttt{forwarding\_unit} detecta
dependencias entre a instrucao no EX e as instrucoes em MEM/WB. Se houver
dependencia, o operando e encaminhado diretamente, evitando stall.

\textbf{Internal forwarding WB$\rightarrow$ID}: Quando uma instrucao no WB
escreve em um registrador que uma instrucao no ID esta lendo no mesmo ciclo,
a escrita sincrona do banco ainda nao ocorreu. Um mux combinacional seleciona
o dado do WB diretamente para o ID/EX register.

\subsection{Load-Use Hazard}
Quando um load esta no EX e a instrucao no ID depende do resultado, a
\texttt{hazard\_unit} gera stall no IF/ID e flush no ID/EX, inserindo uma
bolha no pipeline.

\subsection{Control Hazards}
Branches sao resolvidos no EX. Quando um branch e tomado ou um jump e
executado, flush e aplicado no IF/ID e ID/EX, e o PC e redirecionado.

\section{Sinais de Debug}

\begin{tabularx}{\textwidth}{|l|c|X|}
\hline
\textbf{Sinal} & \textbf{Bits} & \textbf{Descricao} \\
\hline
\texttt{pc\_debug\_o} & 32 & PC atual (estagio IF) \\
\texttt{instr\_debug\_o} & 32 & Instrucao atual (estagio ID) \\
\texttt{alu\_result\_debug\_o} & 32 & Resultado da ALU (estagio EX) \\
\texttt{alu\_carry\_debug\_o} & 1 & Carry da ALU (debug do somador) \\
\texttt{alu\_overflow\_debug\_o} & 1 & Overflow da ALU (debug do somador) \\
\texttt{reg\_debug\_o} & 32 & Valor do registrador selecionado por \texttt{reg\_sel\_i} \\
\texttt{stage\_if\_pc\_o} & 32 & PC no estagio IF \\
\texttt{stage\_id\_pc\_o} & 32 & PC no estagio ID \\
\texttt{stage\_ex\_pc\_o} & 32 & PC no estagio EX \\
\texttt{hazard\_stall\_o} & 1 & Indicador de stall ativo \\
\texttt{hazard\_flush\_o} & 1 & Indicador de flush ativo \\
\hline
\end{tabularx}

\section{Validacao}

\subsection{Testbench Automatizado}
O arquivo \texttt{tb\_riscv\_cpu.v} implementa um testbench com iverilog que
carrega um programa de 49 instrucoes cobrindo todas as instrucoes exigidas,
executa 70 ciclos e verifica 31 registradores e 2 posicoes de memoria.

\subsection{Simulacao no Digital}
O arquivo \texttt{test\_circuit.dig} contem um circuito pronto para o simulador
Digital, com ROM pre-carregada, RAM dual-port, Clock e probes de debug.
Para executar: abrir o arquivo, com \texttt{Reset=1} clicar no Clock uma vez,
mudar \texttt{Reset=0}, e clicar no Clock repetidamente observando os probes.

\subsection{Resultados}
Todos os 32 testes passaram, cobrindo: aritmetica, logica, deslocamento,
memoria, controle (branch/jump), forwarding e load-use hazard.

\subsection{Como Reproduzir}
\begin{lstlisting}[language=bash]
iverilog -g2012 -s tb_riscv_cpu -o tb.out \
  tb_riscv_cpu.v alu.v branch_comparator.v \
  control_unit.v forwarding_unit.v hazard_unit.v \
  instruction_decoder.v pipeline_regs.v \
  program_counter.v register_file.v riscv_cpu.v
vvp tb.out
\end{lstlisting}


\end{document}
